\section{Conclusion}
\label{sec:conclusion}

We have presented GVCCTurbo, a zero-shot generative video codec that retains
the codebook-driven backbone of GVCC while reducing the cost of its expensive
video-generator trajectory. The organizing principle is trajectory
preservation: keep the SDE/codebook correction grid that carries the
compressed residual, but avoid unnecessary generator work where the missing
drift can be reconstructed reproducibly. Two mechanisms instantiate this
principle---unifying multi-resolution stages in a single latent domain, and
caching the trajectory endpoint rather than the velocity for predictive
skipping. Both are simple in implementation, require no retraining, and leave
the bitstream format unchanged beyond the optional DLC1 transition codebook.

The experiments support two broader lessons. First, multi-resolution coding is
viable for generative compression only when all stages share the same latent
target; otherwise the low-resolution stage optimizes a different object from
the full-resolution decoder. Second, model-call skipping should cache the
clean endpoint, not the velocity: the endpoint remains a stable destination,
whereas the velocity is tied to a stale state. These lessons are codec-specific
and do not follow from ordinary fast-sampling recipes.

\paragraph{Limitations and future work.}
The decoder is still substantially slower than learned codecs of the
DCVC family ($\sim 10^2$ vs.\ $\sim 10^{-1}$~s/GOP), and the main
remaining cost is the per-generator-call latency of the high-resolution
stage. Two directions appear most promising. First, the predictive-skip
mechanism's one-step-radius observation rules out simple multi-step
caching, but does not preclude a learned step-distillation that retrains
a few-step generator compatible with the encoder's per-step
reproducibility requirement---this is an open research direction with no
clean fix today. Second, a lightweight generator natively trained at
1080p (the closest existing candidate at the time of writing,
LTX-Video~\cite{ltxvideo2025}, is a strong starting point) would change
the backbone-vs-codec balance of the cost; we plan to investigate this
trade in future work.

We release code, full per-GOP measurements, and codebook bitstreams to
support reproducibility and further investigation along the multi-
resolution and predictive-skip axes.
